\chapter[Considerações Finais]{Considerações Finais}
\label{cap:consideracoes}

\section[Conclusões]{Conclusões}

Conforme exposto na Seção \ref{sec:objetivos}, este trabalho teve como objetivo o desenvolvimento de uma aplicação web capaz de intermediar o contato entre tutores de diferentes áreas de conhecimento e pessoas interessadas em receber tutoria. Por meio dessa aplicação tornou-se possível a busca e conexão entre aprendizes e tutores, obter recomendações de tutores, realizar o agendamento de sessões de tutoria, fornecer de \textit{feedback} mútuo e haveria também a utilização de gamificação para estímulo da motivação e buscar sustentar o engajamento a longo prazo dos usuários.

A pesquisa aplicada possibilitou o aprofundamento dos conhecimentos sobre as áreas relacionadas à aplicação. A utilização da metodologia XP em conjunto com o Kanban permitiu a organização e o acompanhamento das atividades realizadas durante o desenvolvimento.

Durante a fase de desenvolvimento descrita na Seção \ref{sec:desenvolvimento}, cada US foi implementada de maneira a atender os objetivos propostos. Todas as funcionalidades priorizadas foram implementadas, sendo algumas até mesmo melhoradas se comparadas aos critérios de aceitação originais. A aplicação possui sua qualidade atestada pelos testes unitários, de integração e E2E, além da análise estática do código. Além disso, os testes de usabilidade demonstraram que os elementos de gamificação foram bem recebidos e conseguem o efeito desejado de estímulo da motivação do usuário, conforme o relato de uma participante. O \textit{deploy} foi realizado, então a aplicação de fato pode ser utilizada por meio de navegadores de Internet.

Em relação aos objetivos específicos definidos na Seção \ref{sec:}, a implementação dos mecanismos de busca e filtragem de tutores, juntamente com o sistema híbrido de recomendação, possibilitou a localização de tutores e a personalização das recomendações apresentadas aos aprendizes. A implementação dos elementos de gamificação, como conquistas, níveis e títulos, buscou estimular a motivação e o engajamento dos usuários na aplicação. O sistema de agendamento de tutorias, por sua vez, permitiu a organização das disponibilidades dos tutores e o gerenciamento das solicitações de sessões. Por fim, os mecanismos de comunicação e avaliação mútua possibilitaram a interação entre tutores e aprendizes, contribuindo para a construção de uma reputação dentro da plataforma.

Diante dos resultados obtidos, considera-se que o objetivo geral do trabalho foi alcançado, uma vez que a aplicação foi desenvolvida e disponibilizada para utilização por meio de navegadores de Internet, oferecendo os mecanismos necessários para intermediar o contato entre tutores e aprendizes. Quanto à questão de pesquisa definida na Seção \ref{sec:questao}, os resultados demonstram que uma aplicação web pode auxiliar na localização de tutores, na organização de encontros e na disponibilização de estímulos à participação por meio da gamificação.

Embora os elementos de gamificação tenham sido avaliados positivamente durante os testes de usabilidade, incluindo o relato de uma participante que afirmou se sentir reconhecida e incentivada ao utilizar a aplicação, esses resultados não são suficientes para comprovar seus efeitos sobre o engajamento em longo prazo. A comprovação dessa manutenção de engajamento demanda avaliações realizadas durante períodos prolongados de utilização, não contempladas pelos testes deste trabalho. Dessa forma, a questão de pesquisa foi respondida parcialmente, uma vez que a solução tecnológica proposta foi concretizada, mas seus efeitos engajadores em longo prazo ainda demandam investigação.

\section[Trabalhos Futuros]{Trabalhos Futuros}

A partir da aplicação desenvolvida neste trabalho, identificam-se possibilidades de evolução, sendo algumas destacadas como trabalhos futuros interessantes para este projeto:

\begin{itemize}

\item \textbf{Sistema de Pagamento}: Com o objetivo de possibilitar que tutores possam cobrar por suas tutorias, poderia ser implementado um sistema de pagamento na aplicação. Isso envolveria o estudo de diferentes modelos de pagamento, a integração com serviços de processamento de transações financeiras e a definição de mecanismos para assegurar a integridade e a segurança das operações. Além disso, poderiam ser investigadas estratégias de gerenciamento de pagamentos, cancelamentos e reembolsos, assim como aspectos relacionados à proteção dos dados financeiros dos usuários. Além de possibilitar a remuneração dos tutores pelas sessões realizadas, poderia ser investigada a viabilidade de um modelo de negócio baseado na retenção de uma porcentagem dos valores das transações pela aplicação, visando sua sustentabilidade financeira.

\item \textbf{Sistema de Videoconferência}: Com o objetivo de permitir que tutores e aprendizes possam realizar sessões de tutoria à distância diretamente pela aplicação, poderia ser implementado um ambiente integrado de videoconferências na aplicação. Isso envolveria o estudo de tecnologias de comunicação audiovisual em tempo real, bem como a investigação de aspectos relacionados à latência, à estabilidade das conexões, à segurança e à privacidade das comunicações. Adicionalmente, poderiam ser analisados recursos de apoio à aprendizagem colaborativa durante as sessões, como compartilhamento de tela e ferramentas interativas. 

\end{itemize}
